\section{Superscalar Execution}

\begin{definition}
    A $N$-\textbf{Superscalar Processor} is a processor that can issue more than one ($N$) instruction per clock cycle.
\end{definition}

\begin{remark}
    Note that superscalar execution is orthogonal to OoO execution. A processor can be in-order or out-of-order and scalar or superscalar.
\end{remark}

\begin{important}
    The main challenge of superscalar processors is to find independent instructions to execute in parallel. (For example trough Forwarding, Software-Optimizations, Speciulation or Multi-Threading)
\end{important}

\begin{remark}
    The \textbf{tradeoffs} of Superscalar execution are IPC vs Complexity.
\end{remark}